\documentclass[12pt]{article}
\usepackage[utf8]{inputenc}
\usepackage[T1]{fontenc}
\usepackage{hyperref}
\title{Potawatomi Language Resources for K-12 Curriculum Development}
\author{Compiled by Claude Sonnet 5.5 (claude-sonnet-5-5) for Andr\'as Kornai}
\date{}
\begin{document}
\maketitle

\section{Introduction}
Potawatomi (Bod\'ewadmimwen; ISO 639-3 \texttt{pot}, Algic) is an endangered language of the United States, spoken by perhaps a hundred people, mostly elderly, and increasingly learned as a second language. Two tribal language programmes (Citizen Potawatomi Nation and the Pokagon Band) produce most of the teaching material available online. This report is based on a single web search and is not exhaustive.

\section{Primary Educational Resources}
\begin{itemize}
\item Citizen Potawatomi Nation Language Department: Bod\'ewadmimwen offered to Shawnee Public Schools sixth graders. The department has offered a course to Oklahoma public schools since 2017, designed to meet world-language credit standards; it is a light introduction (animals, weather, numbers, everyday vocabulary) using flashcards, matching games, video and audio. \url{https://www.potawatomi.org/blog/2022/01/12/bodewadmimwen-now-offered-to-shawnee-public-schools-sixth-graders/}
\item Pokagon Band, Education Department: ``Want to jump start your language learning?'' (language classes, apprenticeship programme, coloring books, animated videos). \url{https://www.pokagonband-nsn.gov/departments/education/higher-education/learning-the-language/}
\end{itemize}

\section{Linguistic and Reference Materials}
The search summary mentions online dictionaries with audio from both nations (the CPN dictionary and Pokagon's Wiwkw\'ebth\"egen); I did not obtain their exact URLs from the results and so do not list them as resources here. No grammar or orthography document was located in this search.

\section{Religious and Community Resources}
\begin{itemize}
\item Citizen Potawatomi Nation, language-learning tips from the Language Department. \url{https://www.potawatomi.org/blog/2020/11/13/potawatomi-language-learning-tips/}
\item Citizen Potawatomi Nation, page by Language Director Justin Neely. \url{https://potawatomi.org/?p=26476}
\end{itemize}
No scripture translations were found in this search.

\section{Digital and Audio Resources}
The search summary reports a self-paced Moodle course, YouTube channels, Memrise and Quizlet sets from CPN, and the Woolaroo app (Google Arts \& Culture) with 900+ words; specific URLs were not in the results and are not listed. The Pokagon page above also points to online dictionary, video and coloring-book material.

\section{Curriculum Development Context}
Potawatomi has real school presence in Oklahoma (CPN course in public middle schools) and at the Pokagon Band's school partnerships, but available courses are introductory. A full K-12 curriculum would have to be built on the tribal programmes' dictionaries and courses; bilingual literature such as Jim Thunder Sr.'s \emph{Mikw\'end\"ewnen: Memories} was mentioned in the search summary without a link. Contact the language departments directly for current materials.
\end{document}